% Transcription — fonds Alexandre Grothendieck, shelfmark 9, batch 2 (pages 21-28).
% Established from the facsimile archives/batches/9/batch-02.pdf, cut from a
% copy of 9.pdf that was fetched through the project's own relay
% (https://grothendieck-relay.onrender.com/source/9.pdf) rather than from
% grothendieck.umontpellier.fr directly; per the skill transcribe-grothendieck.
% The batch file carries no cover sheet: PDF page k is the archivists' page
% 20+k, and their pencil numbers bottom-left (« 21 » ... « 28 ») agree on
% every leaf. The folder is twenty-eight pages; this is the second and last
% batch. Batch 1 (pages 1-20) is being transcribed separately and did not
% exist when this one was made, so the notation below was fixed from these
% eight pages alone (and from the last line of page 20, read only to know
% where page 21 starts: « d'où σ_A ∈ Γ(det t_A)^{1-p} = Γ ω_A^{(p-1)} »).
%
% Pass: Opus 5.5 (claude-opus-5-5), 2026-09-26 — first pass, unchecked against the pages by a human.
%
% The folder sits in the inventory group "4-9" « Cristaux », dated
% [à partir de 1958- vers 1971]; the folder has its own date, which
% \dating{} carries verbatim.
%
% Pages skipped: none. Pages summarised: none. All eight leaves are his, one
% continuous run in blue-black ink, unpaginated by him. Page 25 is cancelled
% whole by three long diagonal strokes and page 23's closing Lemma by four;
% both are transcribed and the cancellation recorded once. Fast drafting
% throughout: the formulas carry, much connective prose does not and is
% left \ill{}.
%
% What the batch is about. The end of a discussion of the Hasse invariant
% of an abelian scheme A/S in characteristic p and the start of a
% construction of a canonical isomorphism det T_p(A)_ét ≃ det T_p(A*)_ét
% for A ordinary.
%   21  The section σ_A of ω_A^{⊗(p-1)} « mérite le nom d'invariant de
%       Hasse »; its zeros are exactly the s with A_s non-ordinary. From
%       0 → t_A^∨ → DR^1(A/S) → t_{A*} → 0 and det DR^1 ≃ DR^{2g} ≃ O_S,
%       ω_A ≃ ω_{A*}; claimed compatible with the p-structures, sending
%       σ_A to σ_{A*}, and inducing det(pA)_ét ≃ det(pA*)_ét for A ordinary.
%   22  T_p(A)_ét = lim (p^ν A)_ét; the extension 0 → T_p(A)_rad → T_p(A)
%       → T_p(A)_ét → 0 and its Cartier dual, T_p(A)_rad ≃ D(T_p(A*)_ét);
%       wish for a canonical isomorphism.
%   23  det T_p(A)_ét ≃ det T_p(A*)_ét; the analogy with ℓ ≠ p
%       (det T_ℓ(A) ≃ det T_ℓ(A*) ≃ Z_ℓ(g), α_ℓ = β_ℓ(φ)/deg φ for a
%       polarisation φ); definition α_p(φ) = β_p(φ)/√deg φ̃; « À prouver :
%       a) α_p(φ) ne dépend pas de φ »; a cancelled Lemma on det T_p(u)_ét
%       for an endomorphism u, marked « faux » in the margin.
%   24  Proof of a): ψ̃ = φ̃u, u in End(A)⊗Q, u' = φ̃uφ̃^{-1} (Rosati);
%       reduction to Δ(u)/δ(u) = 1, Δ(u)² = δ(u)², hence Δ(u) = ±δ(u); the
%       sign to be settled by positivity (u > 0 ⇒ Δ(u) = +δ(u)).
%   25  (cancelled) an alternative via the p-adic theorem / Weil's
%       characteristic polynomial and a sign character ε on F*.
%   26  Δ(u) polynomial in the coordinates of u with integral values, the
%       positivity argument over R, cqfd; b) α_p(φ) is an isomorphism iff
%       v_p(card Coker T_p(φ̃)_ét) = v_p(√deg φ̃).
%   27  Coker T_p(φ̃)_ét ≃ (Ker φ̃)(p)_ét, autoduality of Ker φ̃, (card)² =
%       deg; the boxed canonical isomorphism α_A : det T_p(A)_ét ⥲
%       det T_p(A*)_ét (*), compatible with base change; for u an
%       automorphism det T_p(u)_ét det T_p(u*)_ét = 1; « Attention » for
%       endomorphisms.
%   28  Functoriality (bracketed), the boxed det T_p(u)_ét det T_p(u*)_ét =
%       deg u; c) α_{A*} = α_A^{-1} to check; d) compatibility with the
%       isomorphism from De Rham duality of page 21 (« à vérifier »).
%
% Notation fixed once for the batch: A* the dual abelian scheme (his star);
% \underline{t}_A, \underline{\mathcal{O}}_S underlined as he writes them;
% \omega_A; \mathrm{DR}^1(A/S); T_p(A)_{\text{ét}}, T_p(A)_{\mathrm{rad}};
% {}_{p}A for his left subscript; \tilde\varphi : A → A* the map of a
% polarisation φ; u' (p. 24) and u* (pp. 24-28) for the Rosati transpose,
% both as written; Δ, δ, ε as on the pages. He writes « det » and « dét »
% indifferently; both are set \det. Greek letters are in math mode.
%
% Readings to check first: « p-structures » on p. 21 (a small stroke above
% the p); the verb after « permet la » on p. 22; « Si ... au lieu de p » on
% p. 23 and the struck bracket after it; the name after « Or par » on p. 24;
% most of p. 25; « comme ... ordinaire » and the bracket « par transport ? »
% on p. 27; the first two lines of p. 28 and its left-margin addition.

\documentclass[11pt,a4paper]{article}
% Le préambule commun à toutes les transcriptions du fonds Grothendieck.
%
% Il tient en une idée : **l'incertitude se compose, elle ne se résout pas.**
% Une transcription qui lisse un mot douteux a détruit la seule chose pour
% laquelle on la faisait — un lecteur doit pouvoir savoir, à chaque endroit,
% ce qui était sur le feuillet et ce qui a été ajouté. Les six macros qui
% suivent sont donc l'appareil critique, et non de la décoration.
%
% Les mêmes commandes sont reconnues par `scripts/render.mjs`, qui produit la
% vue de lecture du site. Une macro ajoutée ici doit l'être aussi là-bas,
% faute de quoi le rendu échouera bruyamment — ce qui est voulu : un rendu
% silencieusement faux est pire qu'un rendu absent.

\ProvidesPackage{grothendieck}[2026/08/08 Transcriptions du fonds Grothendieck]

\usepackage{amsmath,amssymb,amsthm}
\usepackage{mathrsfs}
\usepackage{stmaryrd}
% Les diagrammes commutatifs sont partout dans ce fonds : c'est la langue
% même de la géométrie algébrique de Grothendieck, et une transcription qui
% les rendrait en prose ne transcrirait rien.
\usepackage{tikz-cd}
% Français par défaut — c'est la langue des feuillets — mais l'anglais est
% chargé aussi : la traduction et le résumé sont des documents anglais, et la
% typographie française (espaces avant les deux-points) y serait une faute.
% Ces documents basculent par \selectlanguage{english} en tête de corps.
\usepackage[english,french]{babel}
\usepackage{fontspec}
\usepackage{geometry}
\geometry{a4paper,margin=25mm,marginparwidth=28mm}
\usepackage{marginnote}
\usepackage{xcolor}
% ulem, not soul. `soul`'s \st and \ul are built on a character-by-character
% scan that XeLaTeX's font machinery breaks: the document fails with an
% undefined control sequence rather than degrading. ulem does the same two jobs
% and is Unicode-engine safe. `normalem` stops it hijacking \emph, which the
% transcriptions use for Grothendieck's own emphasis.
\usepackage[normalem]{ulem}
\usepackage{hyperref}
\hypersetup{colorlinks=true,linkcolor=black,urlcolor=black,pdfborder={0 0 0}}

% --- Métadonnées du lot -----------------------------------------------------
% Reprises telles quelles par le script de rendu, qui les affiche en tête de
% la vue de lecture. Elles fixent ce que le fichier prétend couvrir : sans
% elles, une transcription flotte sans point d'attache dans un fonds de seize
% mille feuillets.

\newcommand{\folder}[1]{\gdef\grothFolder{#1}}
\newcommand{\batch}[1]{\gdef\grothBatch{#1}}
\newcommand{\leaves}[2]{\gdef\grothFirst{#1}\gdef\grothLast{#2}}
\newcommand{\foldertitle}[1]{\gdef\grothFoldertitle{#1}}
\gdef\grothFolder{?}\gdef\grothBatch{?}\gdef\grothFirst{?}\gdef\grothLast{?}\gdef\grothFoldertitle{}

% --- Le résumé --------------------------------------------------------------
%
% Il ouvre la lecture modernisée, sous le titre « Résumé », et il est composé
% en petit. Ce n'est pas un ornement : c'est le seul passage du document qui
% ne soit pas des mathématiques, et un lecteur doit pouvoir le reconnaître
% avant de l'avoir lu — puis le sauter s'il connaît déjà le sujet, ou s'y
% arrêter s'il ne le connaît pas. Le filet qui le suit marque la reprise du
% corps.

\newenvironment{resume}
  {\small\setlength{\parskip}{0.45em}}
  {\par\medskip\hrule\medskip}

% --- Mots-clés ---------------------------------------------------------------
%
% En anglais, en clôture du résumé de la lecture modernisée : le vocabulaire
% moderne sous lequel chercher ce que les feuillets construisent. Le manifeste
% du site (`npm run manifest`) les extrait pour étiqueter la cote entière —
% cette ligne est la source unique des tags, il n'y a pas de fichier à côté.
\newcommand{\keywords}[1]{\par\smallskip\noindent
  {\footnotesize\textsc{Keywords} --- \textit{#1}}\par}

% --- L'appareil critique ----------------------------------------------------

% Le numéro de feuillet, en marge. C'est l'ancre qui permet de régler un
% désaccord sur une formule en regardant *un* feuillet plutôt qu'une chemise
% de six cents. Le site s'en sert aussi pour tourner les pages du fac-similé
% quand on fait défiler la transcription.
\newcommand{\leaf}[1]{%
  \marginnote{\footnotesize\color{gray!70}\textsf{#1}}%
  \label{leaf:#1}\ignorespaces}

% Illisible : on ne devine pas. Un mot inventé qui se lit comme les autres est
% le pire résultat possible.
\newcommand{\ill}{\textcolor{red!60!black}{[\,\ldots\,]}}

% Lecture douteuse : le mot est donné, et signalé comme incertain.
\newcommand{\uncertain}[1]{\uline{#1}}

% Ajout de l'éditeur — une lettre manquante, un mot sous-entendu.
\newcommand{\add}[1]{\textcolor{blue!50!black}{[#1]}}

% Biffé par Grothendieck lui-même. Ce qu'il a rayé fait partie du document :
% une rature dit souvent où la pensée a changé de direction.
\newcommand{\struck}[1]{\sout{#1}}

% Note du transcripteur, sur l'état matériel du feuillet.
\newcommand{\note}[1]{\par\smallskip{\small\color{gray!60!black}\textit{[#1]}}\par\smallskip}

% Note marginale de Grothendieck, distincte du corps du texte.
\newcommand{\marginal}[1]{\par\smallskip\noindent\hspace{1em}%
  {\small\color{gray!60!black}#1}\par\smallskip}

% --- Titre ------------------------------------------------------------------

\AtBeginDocument{%
  \begin{center}
    {\large\bfseries Fonds Alexandre Grothendieck --- cote n\textsuperscript{o}~\grothFolder}\\[2pt]
    {\itshape\grothFoldertitle}\\[4pt]
    {\small feuillets \grothFirst--\grothLast\ (lot \grothBatch)}\\[2pt]
    {\footnotesize\color{gray!60!black}Transcription établie d'après le fac-similé numérisé
     par l'Université de Montpellier.\\
     Les lectures incertaines sont soulignées, les illisibles notées [\,\ldots\,],
     les ajouts entre crochets.}
  \end{center}
  \vspace{1em}}

\endinput


\folder{9}
\batch{2}
\pages{21}{28}
\dating{[vers 1969]}
\watermark{Édition de démonstration}
\foldertitle{Groupe de Barsotti-Tate et cristaux (jusqu’en été 1969) : notes manuscrites (s.d.)}

\begin{document}

\page{21}

Cette section mérite le nom d'\emph{invariant de Hasse} de $A$. L'ensemble des
zéros de cette section est exactement l'ensemble des $s$ tels que $A_s$ soit
\emph{non ordinaire}.

Quelles sont les relations entre $\omega_A$ et $\omega_{A^*}$ ? On a, grâce à
la suite exacte
\[
0 \to \underline{t}_A^{\vee} \to \mathrm{DR}^1(A/S) \to \underline{t}_{A^*} \to 0
\]
la relation
\[
\det \underline{t}_A^{\vee} \otimes \det \underline{t}_{A^*} \simeq
\begin{array}{c}
\det \mathrm{DR}^1(A/S) \\
\wr \\
\mathrm{DR}^{2g}(A/S) \simeq \underline{\mathcal{O}}_S
\end{array}
\]
donc on a
\[
\omega_A \cdot \omega_{A^*}^{\vee} \simeq \underline{\mathcal{O}}_S ,
\]
i.e.\ on a un isom.\ can.
\[
\omega_A \simeq \omega_{A^*} .
\]
\emph{Je dis que cet isom.\ est compatible avec les
\uncertain{$p$-structures}, — donc transforme} $\sigma_A$ \uncertain{en}
$\sigma_{A^*}$, et \emph{dans le cas $A$ ordinaire}, induit un isom.
\[
\det({}_{p}A)_{\text{ét}} \simeq \det({}_{p}A^*)_{\text{ét}} .
\]
\note{les soulignements sont les siens ; au-dessus du $p$ de
« $p$-structures » un petit trait que l'on n'a pas su lire}

\page{22}

Il s'impose de regarder également, si $A$ est ordinaire
\[
T_p(A)_{\text{ét}} = \varprojlim \struck{\ill{}}\, ({}_{p^{\nu}}A)_{\text{ét}}
\]
(N.B.\ $T_p(A)$ est une extension
\[
0 \to T_p(A)_{\mathrm{rad}} \to T_p(A) \to T_p(A)_{\text{ét}} \to 0 .
\]
On sait, par dualité de Cartier, que dans la suite exacte correspondante
\[
0 \to T_p(A^*)_{\mathrm{rad}} \to T_p(A^*) \to T_p(A^*)_{\text{ét}} \to 0
\]
les termes extrêmes sont duaux l'un de l'autre, et \struck{plus précisément}
cette dualité est compatible avec filtrations, donc
\[
\begin{aligned}
T_p(A)_{\mathrm{rad}} &\simeq D(T_p(A^*)_{\text{ét}}) \\
T_p(A^*)_{\mathrm{rad}} &\simeq D(T_p(A)_{\text{ét}})
\end{aligned}
\]
donc la connaissance de $T_p(A)_{\text{ét}}$ et $T_p(A^*)_{\text{ét}}$ permet
la \ill{} des \emph{deux} facteurs de $T_p(A)$, donc $T_p(A)$ si la base est
parfaite [en général, il y aurait lieu d'expliciter la structure
d'extension\ldots].

On aimerait encore d'établir un isom.\ can.\ \struck{\uncertain{fonctoriel}}

\page{23}

\[
\det T_p(A)_{\text{ét}} \xrightarrow{\ \sim\ } \det T_p(A^*)
\]

[\struck{Soit} Si \ill{} $\ell \neq p$ \ill{} lieu de $p$, on a bien
\[
\det T_\ell(A) \overset{\alpha_\ell}{\simeq} \det T_\ell(A^*) \simeq
\mathbb{Z}_\ell(g) \quad ]
\]
\note{le crochet fermant est suivi d'un signe barré que l'on n'a pas su lire}

\struck{car $T_\ell(A^*) \simeq \det T_\ell(A^*)^{\vee} \otimes
\mathbb{Z}_\ell(1)$, d'où $\det T_\ell(A^*) \simeq \det T_\ell(A)$}
\note{deux lignes annulées par un zigzag, la seconde en outre entourée ; le
« det » de la première est lui-même biffé, et on lit bien $T_\ell(A^*)$ des
deux côtés}

Si $\varphi$ est une polarisation de $A$, \struck{Choisissons une
polarisation} définissant $\tilde\varphi : A \to A^*$ d'où
$T_\ell(\tilde\varphi) : T_\ell(A) \to T_\ell(A^*)$,
\[
\det T_\ell(\tilde\varphi) : \det T_\ell(A)
\xrightarrow{\ \beta_\ell(\varphi)\ } \det T_\ell(A^*) ,
\]
alors l'isom.\ $\alpha_\ell$ n'est autre que $\beta_\ell(\varphi)/\deg\varphi$.
En s'inspirant de cette formule, on définira,
\[
\beta_p(\varphi) = \det\bigl(T_p(\tilde\varphi)_{\text{ét}}\bigr) :
\det T_p(A)_{\text{ét}} \to \det T_p(A^*)_{\text{ét}}
\]
et on considère
\[
\alpha_p(\varphi) = \beta_p(\varphi)\big/\sqrt{\deg\tilde\varphi} :
\det T_p(A)_{\text{ét}} \otimes \mathbb{Q}_p \to
\det T_p(A^*)_{\text{ét}} \otimes \mathbb{Q}_\ell
\]
\note{les deux « $\otimes\mathbb{Q}$ » sont ajoutés après coup, le premier
au-dessus de la flèche, le second sous la ligne, entouré et relié ;
l'indice $\ell$ du second est tel qu'écrit}

À prouver : a) $\alpha_p(\varphi)$ ne dépend pas de $\varphi$.

\struck{\ill{} formellement du \ill{}}

\textbf{Lemme}. Soit $u$ un \struck{automorphisme} \struck{\ill{}} endom.\ de
$A$, d'où
\[
\Bigl[\det T_p(u)_{\text{ét}} : \det T_p(A)_{\text{ét}} \to
\det T_p(A)_{\text{ét}}\Bigr] \in \mathbb{Z}_p .
\]
On a
\[
\left\{
\begin{array}{l}
\deg u \text{ est un carré parfait} \\
\det\bigl(T_p(u)_{\text{ét}}\bigr) = \sqrt{\deg u}
\end{array}
\right.
\]
i.e.
\[
\deg u = \bigl(\det(T_p(u)_{\text{ét}})\bigr)^2 \quad \text{et} \quad
\det T_p(u)_{\text{ét}} \text{ est un entier} \geq 0 .
\]
\marginal{faux}
\note{le lemme entier, depuis « Soit $u$ », est annulé par quatre longs traits
obliques ; « faux », souligné, est écrit dans la marge gauche, à hauteur de
l'accolade}

\page{24}

En effet, soit $\psi$ une autre polarisation, donc
$\tilde\psi = \tilde\varphi \circ u$, où
$u \in \operatorname{End}(A) \otimes_{\mathbb{Z}} \mathbb{Q}$, \struck{\ill{}}
et la relation $\tilde\psi' = -\tilde\psi$ (pour : $\tilde\varphi' =
-\tilde\varphi$) donne
\[
\begin{array}{c}
u'\tilde\varphi' = -\tilde\varphi u \\
\| \\
-u'\tilde\varphi
\end{array}
\qquad \text{i.e.} \qquad u' = \tilde\varphi u \tilde\varphi^{-1} .
\]
On veut prouver [posons $\delta(\tilde\varphi) = \sqrt{\deg\tilde\varphi}$,
$\alpha_p(\tilde\varphi) = \Delta(\tilde\varphi)$]
\[
\begin{array}{l}
\Delta(\tilde\psi)/\delta(\tilde\psi) = \Delta(\tilde\varphi)/\delta(\tilde\varphi) \\
\| \\
\Delta(\tilde\varphi)\cdot\Delta(u)/\delta(\tilde\varphi)\,\delta(u)
\end{array}
\]
qui équivaut :
\[
\Delta(u)/\delta(u) = 1 .
\]
Or par \ill{} on a
\[
\Delta(u)\,\Delta(u^*) = \delta(u)^2 = \deg u
\]
\note{le « $= \deg u$ » est écrit sous $\delta(u)^2$, relié par un signe
d'égalité vertical ; un $\Delta$ isolé, repassé, précède la formule}

d'autre part, comme $u^* = \tilde\varphi u \tilde\varphi^{-1}$, on a ici
$\Delta(u^*) = \Delta(u)$, donc \struck{\ill{}}
\[
\Delta(u)^2 = \delta(u)^2
\]
d'où $\Delta(u) = \pm\delta(u)$. Il faut encore \uncertain{exorciser} le signe,
il faut \struck{prendre} dans un
$\operatorname{End}(A) \otimes_{\mathbb{Z}} \mathbb{Q}$, muni de l'involution
$u \mapsto u' = \tilde\varphi u \tilde\varphi^{-1}$, montrer que $u > 0$ (i.e.\
\struck{\ill{}} $u = \sum v_i v_i'$) \struck{$u = u'$} implique
\[
\Delta(u) = +\delta(u)
\]
\note{« $+\delta(u)$ » est souligné deux fois}
[N.B.\ $\delta(u)$ est bien défini, car $\deg u$ est bien un carré parfait si
$u = u'$]

\struck{et l'identité $\Delta(u)^2 = \delta(u)$ est une identité des}

\page{25}

\note{toute la page est annulée par trois longs traits obliques ; on la
transcrit telle qu'elle se lit}

\struck{Comme $T_p(u)$ et $T_p(u^*)$ sont} de \uncertain{même}

En effet, d'après le théorème $p$-adique ou th.\ caractéristique de Weil, on
aurait
\[
(*) \qquad \deg u = \det T_p(u)_{\text{ét}}\, \det T_p(u^*)_{\text{ét}} ,
\]
et il reste à prouver que
\begin{enumerate}
\item[1°)] $T_p(u)_{\text{ét}} = T_p(u^*)_{\text{ét}}$
\item[2°)] $\varepsilon(u) = 1$
\end{enumerate}
\marginal{ce qui implique
$\det(T_p(u)) = \varepsilon(u)\sqrt{\deg u}$, \quad $\varepsilon(u) = \pm 1$}
\note{ajout à droite de 1°), relié par une accolade ; « de
\uncertain{même} », en haut à droite, est relié par un trait après
« $p$-adique »}

D'ailleurs on peut se borner à \ill{} $u \in \operatorname{Hom}(A,A)
\otimes_{\mathbb{Z}} \mathbb{Q} = F$, \note{un $\varepsilon$ est écrit sous le
signe d'égalité} et alors $\varepsilon$ est un \ill{}
\[
F^* \to (+1, -1) ,
\]
\struck{donc} [\uncertain{position} des \ill{}] \ill{} fonction \ill{}
\struck{qui est une \ill{} \ill{} tel \ill{}} \struck{\ill{}} deux fonctions,
\note{« se borner à » est ajouté au-dessus de la ligne, relié par un trait ; les deux
dernières lignes, déjà annulées par les traits obliques, sont en outre biffées
et raturées d'un zigzag}

\page{26}

Or pour $u$ variable, $\Delta(u)$ est un polynôme dans
\struck{$\mathbb{Q}$}$\mathbb{Z}_p[T]$ ($T = (T_1, \ldots, T_N)$, $N =$
\struck{\ill{}} ${}_{\mathbb{Z}}\operatorname{End}(A)_{\ldots}$)
\note{le mot biffé devant l'indice $\mathbb{Z}$ n'a pas été lu (on attend un
rang) ; un petit mot illisible est écrit en indice après
$\operatorname{End}(A)$}
qui, \ill{}, prend des valeurs entières pour des valeurs entières des
variables. Donc $\Delta(u)$ est un polynôme à coeff.\ \struck{entiers} dans
$\mathbb{Q}$ (et même dans $\mathbb{Q}_p$),
\struck{\ill{} \ill{} à coeff.}
\struck{\ill{} $\mathbb{Z}_{(p\mathbb{Z})}$ [car $\mathbb{Z}_p[T] \cap
\mathbb{Q}[T] = \mathbb{Z}[T]$,} \struck{car $\mathbb{Z}_p \cap \mathbb{Q} =
\mathbb{Z}_{(p\mathbb{Z})}$]}.
\note{ces deux lignes sont biffées et raturées de zigzags}
Il faut prouver que ses valeurs sur le cône des $u > 0$ sont $\geq 0$, et
\struck{pour que} il suffit de le prouver en passant de $\mathbb{Q}$ à
$\mathbb{R}$. Mais dans $\mathbb{R}$ tout élément $\geq 0$ est un carré, d'un
élém.\ \ill{} $\geq 0$, donc
\[
\Delta(u) = \Delta(v^2) = \Delta(v)^2 \in \mathbb{R}^2 = \mathbb{R}^+ ,
\]
cqfd.

b) $\alpha_p(\varphi)$ induit bien un isom.\ \struck{\ill{}} de
$\det T_p(A)_{\text{ét}}$ avec $\det T_p(A^*)_{\text{ét}}$. Cela
\struck{\ill{} que les deux} équivaut : la relation
\struck{(i) $\alpha_p(\varphi)$ applique $\det T_p(A)_{\text{ét}}$ dans
$\det T_p(A^*)_{\text{ét}}$]}
\[
\text{i.e.}\quad
v_p\Bigl(\operatorname{card}\bigl(\operatorname{Coker}
T_p(\tilde\varphi)_{\text{ét}}\bigr)\Bigr) = v_p\bigl(\sqrt{\deg\tilde\varphi}\bigr) .
\]
\note{avant ce « i.e. », un $\beta_p(\varphi)$ est écrit puis raturé ; dans le
conoyau, un premier symbole est noirci et remplacé par $T$}
On a évidemment

\page{27}

\[
\operatorname{Coker} T_p(\tilde\varphi)_{\text{ét}} \simeq
\bigl(\operatorname{Ker} \tilde\varphi\bigr)(p)_{\text{ét}}
\]
Or on sait que $\operatorname{Ker}\tilde\varphi$ est autodual, donc
$(\operatorname{Ker}\tilde\varphi)(p)_{\text{ét}}$ est dual de
$(\operatorname{Ker}\tilde\varphi(p))_{\mathrm{rad}}$, donc comme \ill{}
ordinaire, d'où la relation
\[
\begin{array}{c}
\bigl(\operatorname{card} (\operatorname{Ker}\tilde\varphi)(p)_{\text{ét}}\bigr)^2
= \operatorname{card} \operatorname{Ker}(\tilde\varphi(p)) \\
\hspace{7em}\| \\
\hspace{7em}(\deg\tilde\varphi)(p)
\end{array}
\]
d'où la relation envisagée\ldots\ Cela définit bien un isom.\ canonique
\[
(*) \qquad \boxed{\det T_p(A)_{\text{ét}} \xrightarrow[\alpha_A]{\ \sim\ }
\det T_p(A^*)_{\text{ét}}}
\]
compatible avec changement de base, et isom.\ en $A$.

[Si $u$ est un automorphisme de $A$, et $\check u = u^{*-1}$ \ill{}, \ill{}
on aura \ldots\ $\alpha$, $\det T_p(u)_{\text{ét}} = \det T_p(\check
u)_{\text{ét}}{}^{\alpha}$, donc en \uncertain{identifiant} \ill{} \ill{} à
$\mathbb{Z}_p$ \ill{}, on aura
\[
\begin{array}{c}
\det T_p(u)_{\text{ét}} = \det T_p(u^{*-1})_{\text{ét}} \\
\hspace{7em}\| \\
\hspace{7em}\bigl(\det T_p(u^*)\bigr)^{-1}
\end{array}
\]
donc
\[
\boxed{\det T_p(u)_{\text{ét}}\,\det T_p(u^*)_{\text{ét}} = 1}
\qquad \text{\emph{$u$ un automorphisme de $A$}} \quad ]
\]
\note{au-dessus du second facteur, un « det » est récrit en correction}

Attention, si $u$ est un endomorphisme de $A$, on aura en général
\[
\det T_p(u) \neq \det T_p(u^*)
\]
\note{dans la marge gauche, en regard, un grand signe en forme de « 2 » ou de
« Z »}

\page{28}

\ill{} si $A$ \ill{} \ill{} \ill{} corps alg.\ clos ; donc l'isom.\ \ill{}
\emph{bien} fonctoriel pour les \ill{} de $=$ dimension [dans la mesure où ça
pourrait avoir un sens, les deux foncteurs n'étant pas de $=$ variance en
$A$\ldots]. Cependant, \struck{\ill{}}
\[
\boxed{\det T_p(u)_{\text{ét}}\, \det T_p(u^*)_{\text{ét}} = \deg u}
\]
\marginal{vaut aussi \ill{} \struck{\ill{}} pour \ill{} de $=$ dimension
\ill{} Sch.\,A.\ de}
\note{ajout de la marge gauche, écrit en oblique et relié au texte par un
trait ; lecture très incertaine}

c) Il faudrait vérifier que
\[
\alpha_{A^*} = \alpha_A^{-1}
\]
[\ill{} signe (?)\ldots]. Noter aussi la compatibilité par \ill{} de
\uncertain{$\mathcal{V}A$}.\ ---

d) L'isom.\ $({}_{p}A)_{\text{ét}} \to ({}_{p}A^*)_{\text{ét}}$ induit par
$\alpha_A$ est bien celui déduit plus haut de la dualité en cohomologie de De
Rham (à vérifier\ldots).
\note{« plus haut » renvoie à la page~21, où l'isomorphisme est écrit entre
les déterminants $\det({}_{p}A)_{\text{ét}}$ et $\det({}_{p}A^*)_{\text{ét}}$}

\end{document}
